\documentclass[11pt]{article}
\usepackage[margin=1in]{geometry}
\usepackage{amsmath,amssymb,amsthm}
\usepackage{url}
\usepackage[hidelinks]{hyperref}
\newtheorem{theorem}{Theorem}
\newtheorem{lemma}[theorem]{Lemma}
\newtheorem{corollary}[theorem]{Corollary}
\theoremstyle{definition}
\newtheorem{definition}[theorem]{Definition}
\newtheorem{remark}[theorem]{Remark}
\newcommand{\F}{\mathbb{F}_2}
\newcommand{\Deg}{\mathrm{Deg}}
\title{Grok 4.7 hourly work log --- 04:55 KST\\
\large Review of \texttt{303d735}..\texttt{186fe26} and a polynomial-calculus degree transfer\\ through the pinned E/B compiler (the GF(2)/Gaussian-elimination lane)}
\author{Grok 4.7 (automated reviewer)}
\date{2026-11-28 KST (2026-11-27 PST)}
\begin{document}
\sloppy
\maketitle

\noindent\textbf{Executive status.} P vs NP remains \textbf{OPEN}. Nothing here claims P$\neq$NP, P$=$NP, a Clay result, or an unrestricted lower bound. The new theorem concerns \emph{restricted} proof systems: polynomial calculus (PC) over a field, and the hybrid ``bounded factor + Gaussian elimination + shallow affine branching'' calculi that the repository's lanes use. Root \texttt{main.tex} was \emph{not} edited. Labels follow the current AGENTS.md (\texttt{27f2f5e}): KNOWN THEOREM, PROVED WITH SCOPE, COMPUTATIONAL, CONDITIONAL, CANDIDATE, OPEN, FALSIFIED.

\section{Review log}
\subsection*{Commits seen}
The previous Grok work log is \texttt{bab8a1d2566} (2026-11-27 11:05:37 PST). It reviewed through \texttt{b5812b4} and deferred the transport track. This run covers \textbf{125 non-Grok commits} on \texttt{main}:
\begin{itemize}
\item 19 transport-track commits \texttt{303d735}..\texttt{3f5d8a4} (10:57--11:03 PST; 18 ``Add \ldots'' artifacts plus report \texttt{f002c37});
\item 106 commits \texttt{307efbd}..\texttt{186fe26} (11:06--11:51 PST): T1 36, T2 24, T3 21, T4 11, X01-33 9, C04-28 4, and 1 root-level \texttt{DYNAMIC-COMPATIBILITY-\allowbreak OBLIGATION1-\allowbreak SOFT-\allowbreak ATTENTION-\allowbreak RESIDUAL1.md} (\texttt{af96de2}).
\end{itemize}
Reports read (in full or in large part): T3 transport \texttt{f002c37} (\path{research/agents/ground-up-unsat-bounds/2026-11-27_T3_native_transport_replay/report.md}); T2/C08-34 orientation \texttt{de02d5d}; T1 FP51 energy gap \texttt{fde204c}; X01-33 \texttt{34c7bf8}; C04-28 all-split refusal \texttt{d862b92}; T3 batched probe \texttt{8f51f25}; T1 C07-27 multifiber \texttt{d21459d}; T2/C05 affine orbit \texttt{22c93cf}; T1/L02 NAE sewing \texttt{21d3f8c}; T3/L05-28 paired projection \texttt{eeec07b}; T4 restart credit \texttt{9dca2f3}; root DYNAMIC file \texttt{af96de2}. The remaining commits are artifacts (certificates, results, checkers). For those I checked commit messages only. AGENTS.md is unchanged since \texttt{27f2f5e}: its first 19\,KB are byte-identical to last run's copy, and the connector truncated the tail. Commits landing after \texttt{186fe26} (from 11:53 PST, e.g.\ \texttt{0049ad6}, \texttt{b85f530}, ``Record T1 guarded-pool BLOCK \ldots'') are left for the next run. Earlier Grok files are intact (blobs \texttt{1a98146}, \texttt{086fad3}).

\subsection*{New findings}
\begin{enumerate}
\item \textbf{T3 transport theorem (\texttt{f002c37}): math correct, label non-compliant.} The inequality $(s{+}O(1))\,E[T_{\rm local}]+(a{+}O(1))\,E[\sum_j\Delta_j]\ge Dh-O(D)$ follows from the stated per-cut simulation. Distinct diagonal transcripts come from the rectangle property, and prefix-free coding with Kraft gives $E|\text{transcript}|\ge\log_2 r!$. With no remote access, it yields the stated $\Omega(N^2)$. This is the classical Hennie (1965) crossing-sequence bound, phrased in Yao's communication model, and the report cites Hennie itself. The status \emph{``PROVED/EXACT''} is not an AGENTS.md label. It should read KNOWN THEOREM (instantiated) or PROVED WITH SCOPE with a review status. One unchecked hypothesis: the simulation needs the tape content between monitored cuts to be input-independent or owned by one side. I did not read manifest \texttt{229ca31} to confirm this.
\item \textbf{Other label/format issues.} T3 batched probe (\texttt{8f51f25}) uses a bare ``PROVED:'' with no ``WITH SCOPE'' and no review status. Its $k$-lane bit-mask propagation formula is correct; it is ordinary bit-parallel unit propagation. T1 C07-27 (\texttt{d21459d}) uses a bare ``EXACT'' in the body. The root-level DYNAMIC file (\texttt{af96de2}) carries no AGENTS label for its main claim (``collapses''; only OPEN appears), is not timestamped, and is not in a role folder.
\item \textbf{Spot-checks (correct).} T1/L02: $\exists p\,[\mathrm{NAE}(A,p)\wedge\mathrm{NAE}(B,\neg p)]\iff\mathrm{NAE}(AB)$, checked exhaustively for $|A|,|B|\le3$ (196 cases). This is the same NAE-splitting identity as the RT01 NAE$_4\to$NAE$_3$ split. The single-private-port rule $p=2-\text{sum}$ is correct. T3/L05-28: if the forbidden set of $\exists v(F_1\wedge F_2)$ is a subcube, one blocking clause is exact. This is elementary. T4 restart (\texttt{9dca2f3}): the ``novelty invariant'' (a UP-closed path's conflict core is not already in $D$) is correct, and it is the standard asserting-clause fact of CDCL. The fixed-$k$ credit $\sum_i2^i\binom ri$ matches known bounded-width learning analyses (Atserias--Fichte--Thurley, JAIR 40 (2011)). The AND-chain falsifier, where the minimal core has $h$ decisions, is correct.
\item \textbf{Shared affine lanes.} T2/C05 (orbit resolution, basis shear), T3 (parity GE; native rank 395/643 and 369/593 consistent), and C04-28/T1 (linear splits) all leave open whether affine/GF(2) reasoning plus shallow branching escapes the stagnation family of the 04:04 KST run. \S2 shows that it does not, for every such calculus whose moves have bounded affine span.
\end{enumerate}

\section{Gap tackled: degree transfer (the named OPEN lemma of the 04:04 KST run)}
\subsection*{Setting}
Notation is as in \path{comments/grok-4.7/2026-11-28/2026-11-28_0404_GROK_WORK.tex}. $F$ is a CNF with clauses of 1--3 literals. $N(F)$ is the normalized native factor system of pinned \texttt{compile()} (\texttt{21b9d53}, blob \texttt{7cb0303}): one carrier per signed component, and one ``exactly two of four signed ports'' factor per B vertex. $N_g:=N\wedge(g=0)$. By the 04:04 locality lemma, carriers are $x_v,z_i,w_{i,0},w_{i,1},f$ ($n+3m+1$ of them), every B factor of gadget $i$ lies on carriers in $\mathrm{vars}(C_i)\cup\{z_i,w_{i,0},w_{i,1},f\}$, and with $f=0$ the canonical completion $\sigma$ satisfies gadget $i$ iff $C_i$ is true. Each $\sigma_y$ depends only on $S(y)$, with $|S(y)|\le3$.

\textbf{PC.} Fix a field $K$. A PC refutation of a polynomial set $A$ is a sequence of polynomials. Each line is an axiom (from $A$ or $x^2-x$), a $K$-linear combination of earlier lines, or $x\cdot q$ for an earlier line $q$. The sequence ends with $1$. $\Deg_K(A)$ is the least possible maximum line degree. Encodings: a clause $C$ becomes $P_C=\prod_{\ell\in C}(1-\hat\ell)$ (zero iff $C$ is true), and a factor becomes its forbidden-word indicators $\delta_a=\prod_{y}(y\text{ or }1-y)$. These are exactly the clause polynomials of the full-word CNF expansion. The gauge is the polynomial $g$.

\begin{lemma}[Local completeness; KNOWN THEOREM (standard), proof included]\label{lem:lc}
Let $A$ be polynomials on a variable set $U$ with $|U|=m$, and let $q$ be a polynomial on $U$ of degree $e$ that vanishes on every Boolean common zero of $A$. Then $q$ has a PC derivation from $A$ of degree $\le\max(e,\,m+\max_{r\in A}\deg r)$.
\end{lemma}
\begin{proof}
$q-\mathrm{ml}(q)$ (multilinearization) is derived from Boolean axioms in degree $\le e$, and $\mathrm{ml}(q)=\sum_a q(a)\delta_a$. If $q(a)\ne0$, pick $r\in A$ with $r(a)\neq0$. Multiplying $r$ by the $\le m$ affine factors of $\delta_a$ and reducing by $x\delta_a\equiv a_x\delta_a$ gives $r(a)\delta_a$ in degree $\le m+\deg r$.
\end{proof}

\begin{theorem}[PC degree transfer through the pinned compiler; PROVED WITH SCOPE, author proof, no independent review]\label{thm:deg}
For every field $K$, every CNF $F$ with clauses of 1--3 literals, and every gauge carrier $g$,
\[\Deg_K(F)\ \le\ \max\{\,3(\Deg_K(N_g(F))+1),\ 12\,\}.\]
Equivalently, if $\Deg_K(F)>12$ then $\Deg_K(N_g(F))\ge\Deg_K(F)/3-1$.
\end{theorem}
\begin{proof}
\emph{Gauge.} $N$ is invariant under $\tau:y\mapsto1-y$ on all carriers. Flipping a carrier flips all of its ports, since port signs are fixed, and weight $w$ becomes $4-w$. So each B factor's forbidden set is closed under complement; (B) below verifies this computationally. Let $\pi$ be a degree-$d$ refutation of $N\cup\{g\}$. Multiply every line by $(1-g)$. Axioms of $N$ become derivable in degree $+1$. The axiom $g$ becomes $-(g^2-g)$. Linear combinations and multiplications commute with the factor $(1-g)$. The result derives $1-g$ from $N$ in degree $d+1$. Applying $\tau$, a degree-preserving affine substitution that maps the axioms of $N$ and the Boolean axioms to themselves, derives $g$. Adding the two gives $1$. So $\Deg(N)\le\Deg(N_g)+1$.

\emph{Substitution.} Let $p_y\in\mathbb Z[x]$ be the multilinear extension of $\sigma_y$ on $S(y)$, read in $K$; then $\deg p_y\le3$ and $p_f=0$. Let $\pi$ refute $N$ in degree $D$, and put $\rho(q):=q(p)$. By induction, every $\rho(q)$, $q\in\pi$, is derivable from $F$ in degree $\le\max(3D,12)$:
(a) For a B-factor axiom $\delta_a$ of gadget $i$, $\rho(\delta_a)$ is a polynomial on $\mathrm{vars}(C_i)$ ($\le3$ variables) of degree $\le12$. It vanishes wherever $C_i$ is true, because $\sigma$ then satisfies the factor. Lemma~\ref{lem:lc} with $A=\{P_{C_i}\}$ gives degree $\le\max(12,6)$. The padded clause $[a,b,b]$ differs from $P_C$ by a Boolean-ideal term of degree $\le3$.
(b) For a Boolean axiom, $p_y^2-p_y$ vanishes on all Boolean points of $\le3$ variables, so it is derivable in degree $6$.
(c) $\rho$ is linear.
(d) For a multiplication $q\mapsto yq$, write $\rho(yq)=\sum_M c_M M\rho(q)$ with $|M|\le3$. This takes $\le3$ variable multiplications from $\rho(q)$, so the degree is $\le3\deg q+3\le3D$.
$\rho(1)=1$, so $\Deg(F)\le\max(3D,12)$.
\end{proof}

\begin{definition}[Affine factor calculus $\mathrm{AFC}_K(k,d,s)$]
\emph{$K=\F$:} a \emph{factor} is a tuple of $\le k$ affine $\F$-forms in the carriers together with an allowed set $R$; equations of any support are factors on one form. A \emph{move} adds any factor implied (on Boolean carrier assignments) by current factors when the linear parts of all forms involved span dimension $\le k$. This includes Gaussian elimination on equations of unbounded support (span 2 per step), bounded-width joins and projections, and T2-style factors on sheared/orbit coordinates when these are affine. The calculus may also branch on affine forms along paths of depth $\le d$, and learn nogoods on $\le s$ affine literals derived by refuting under them.
\emph{Any $K$:} factors on $\le k$ carriers, $K$-linear equations of any support with linear combination, conversion between equations and factors on $\le k$ carriers, and branching and nogoods on carrier literals only.
\end{definition}

\begin{lemma}[Simulation; PROVED WITH SCOPE, author proof]\label{lem:sim}
If $\mathrm{AFC}_{\F}(k,d,s)$ refutes $N_g$, then $\Deg_{\F}(N_g)\le2k+d+s+1$. If $\mathrm{AFC}_K(k,d,s)$ refutes $N_g$, then $\Deg_K(N_g)\le k+d+s+1$.
\end{lemma}
\begin{proof}
($\F$) Fix a linear basis $b_1..b_r$, $r\le k$, of the span used by a move. Every form is affine in $b$, and every factor indicator is a polynomial of degree $\le k$ in $x$. In characteristic 2, $b_j^2-b_j=\sum_{i}(x_i^2-x_i)$, so the Boolean axioms for $b$ are available, and multiplying them by $b$-monomials of degree $\le k-2$ stays in $x$-degree $\le k$. As in Lemma~\ref{lem:lc} with the roles of variables played by $b$, each target indicator equals $\sum_\beta\delta_\beta(b)$ over $b$-words $\beta$ it covers. The forms have independent linear parts, so every $\beta$ is realized and semantic implication over carriers equals implication over free $b$. Each $\beta$ is forbidden by some premise polynomial $P$ of degree $\le k$, and $\delta_\beta P\equiv P(\beta)\delta_\beta$ costs degree $\le2k$. An affine branch $\lambda:=\ell-e$ satisfies $\lambda^2\equiv\lambda$ (char 2). Multiply the refutation of child $\lambda=0$ by $(1-\lambda)$ and that of child $\lambda=1$ by $\lambda$. This yields $1-\lambda$ and $\lambda$, at $+1$ degree per level. A nogood costs $+s$ in the same way.
(Any $K$) As in the 04:04 simulation lemma, factor moves cost degree $\le k$ and equations are degree-1 lines. Equation$\to$factor uses $\delta_a(L-b)\equiv(L(a)-b)\delta_a$ with $L(a)\neq b$, in degree $\le k+1$. Factor$\to$equation interpolates $L-b=\sum_a(L(a)-b)\delta_a$, in degree $\le k$. Literal branching uses $y^2-y$.
\end{proof}

\begin{corollary}[GF(2) and $\mathbb Q$ affine lanes stall on the pinned-compiler random family; PROVED WITH SCOPE relative to KNOWN THEOREMS]\label{cor:stall}
Fix $\Delta\ge5.2$, and let $F$ be a uniformly random 3-CNF with $n$ variables and $\Delta n$ clauses. With probability $1-o(1)$:
(1) $F$ and $N(F)$ are UNSAT (first moment, as at 04:04).
(2) $\Deg_K(F)\ge\varepsilon n$ for every field $K$ (KNOWN THEOREM: Ben-Sasson--Impagliazzo, Comput.\ Complexity 19 (2010), char $\neq2$; Alekhnovich--Razborov, FOCS 2001 / Proc.\ Steklov Inst.\ 242 (2003), char 2, which proves the BSI conjecture).
(3) For every gauge $g$, $\Deg_K(N_g)\ge\varepsilon n/3-1$, so no $\mathrm{AFC}_{\F}(k,d,s)$ with $2k+d+s\le\varepsilon n/3-3$ and no $\mathrm{AFC}_{\mathbb Q}(k,d,s)$ with $k+d+s\le\varepsilon n/3-3$ refutes $N_g$. Since $N_{\rm bits}=\mathrm{poly}(n)$, this includes every budget $c\log N_{\rm bits}$.
(4) Every PC$_K$ refutation of $N_g$ has $\exp(\Omega(n))$ monomials, by the size--degree relation of Impagliazzo--Pudl\'ak--Sgall (Comput.\ Complexity 8 (1999)) with $n+3\Delta n+1=\Theta(n)$ carriers and initial degree 4 (KNOWN THEOREM applied).
\end{corollary}

\begin{remark}[Pure $\mathbb Q$-GE is vacuous; PROVED WITH SCOPE, trivial]
For \emph{every} STRICT-EB4 instance, the point $y\equiv\tfrac12$ satisfies all signed-edge equations ($y_b=y_a$ or $1-y_a$), all E equalities, and all B sums ($4\cdot\tfrac12=2$). So rational Gaussian elimination without Boolean axioms never refutes, and any $\mathbb Q$-affine refutation needs degree $\ge2$ reasoning. The corollary bounds that reasoning on the random family.
\end{remark}

\subsection*{Computation (COMPUTATIONAL; \texttt{code/degree\_transfer.py}, seed 2026112805, pinned compiler run via the existing \texttt{code/compile\_batch.js} and \texttt{code/compile\_pinned\_21b9d53.js})}
\begin{itemize}
\item[(A)] \textbf{Substitution identity, exact over $\mathbb Z$.} For 30 random 3-CNFs ($n=4..9$, $m\in[2n,6n]$), there are 3160 B factors and 31\,600 forbidden words. For every one, $\mathrm{ml}(\rho(\delta_a))=c\cdot P_{C_i}$ with $c\in\{0,1\}$, and $c=1$ for exactly one word per factor. Also $\max\deg p_y=3$, and the unreduced degree is $\le12$. All pass. Since the identity holds over $\mathbb Z$, it holds over every field.
\item[(B)] Complement closure of each factor's forbidden set: 3160/3160 pass.
\item[(C)] Half point satisfies the $\mathbb Q$-linear system: 30/30 and 20/20.
\item[(D)] \textbf{GF(2) parity relaxation} (edge XOR, E equality, B even parity), exact GE on 20 compiled UNSAT random 3-CNFs ($n=8..14$, $m=6n$, up to $\sim$32\,600 ports): \emph{consistent in 20/20}. Pure GF(2)-GE never refuted. This illustrates the corollary only, since at these $n$ the degree bound is not informative.
\item[(E)] T1/L02 NAE sewing identity: 196/196.
\end{itemize}
Raw output is in \texttt{code/degree\_transfer.out.json}; the run takes about 5\,s.

\subsection*{Limits and what remains OPEN}
\begin{itemize}
\item These are lower bounds for \emph{restricted} systems: PC over a field and bounded-span affine factor calculi with shallow affine branching. They are \textbf{not} lower bounds for algorithms in general, for Frege, for extended resolution, or for P vs NP. The stagnation family is random 3-CNF through the pinned compiler, which is not the two pinned corpus instances ($n=24,40$, where the bound is uninformative).
\item Not covered: \emph{nonlinear} defined coordinates (extension variables), factors over affine forms of span $>k$, DAG-like (non-tree) affine branching (full Res($\oplus$)), and inequality reasoning (matching/Hall, LP/cardinality, flow). Clauses longer than 3 (split chains) need a separate locality argument.
\item The constants ($3$, $12$, $2k$) are not optimized. The KNOWN THEOREMS are cited from the literature, not re-proved. Only the transfer, the simulation, and the gauge step are proved here.
\item \textbf{Next lemma (OPEN).} The same substitution $\rho$ maps squares to squares, so it should transfer Sherali--Adams/SOS degree, with a target of Grigoriev--Schoenebeck linear SOS degree for random 3-CNF. That would cover the inequality lanes (T1 cardinality/Hall/flow, C04-28 matching relaxations) that PC does not capture.
\end{itemize}

\section*{Files}
This file; \texttt{code/degree\_transfer.py}; \texttt{code/degree\_transfer.out.json}. It reuses the existing \texttt{code/compile\_pinned\_21b9d53.js} (blob \texttt{7cb0303}) and \texttt{code/compile\_batch.js} from the 04:04 run.
\end{document}
